\section{Scope and Imported Finite Structure}

\subsection{Purpose of the paper}

This paper gives a finite theorem of Thalean relational saturation on
the registered mechanics associated with
\[
\Gfifteen\cong L(\Petersen).
\]

The result combines three exact finite structures:

\begin{enumerate}[label=(\roman*)]

\item a native carrier language whose lawful completion family
contracts as
\[
8\longrightarrow2\longrightarrow1;
\]

\item a retained five-address-hidden sector selected exactly by the
context-activated native designation response;

\item the normalized Thalean quadratic, which independently contracts
centered registered contrast toward common mode.

\end{enumerate}

The paper does not introduce a new graph or alter the established
native carrier law.

It also does not identify native loading with iteration of the
Thalean quadratic.

That identification was an earlier conjectural possibility.  The
exact finite result is different and stronger: native selection and
quadratic contrast concentration are distinct certified operations
joined through a common finite response algebra.

\subsection{Imported theorem object}

We import the following finite results from the established
AT4val[60,6] theorem chain
\citep{cave2026dodecahedral,cave2026quantumrelativity}.

\begin{enumerate}[label=(\roman*)]

\item The quotient core is
\[
\Gfifteen\cong L(\Petersen).
\]

\item The transport-induced sector--edge incidence matrix
\[
M\in\{0,1\}^{15\times30}
\]
has row weight \(14\) and column weight \(7\).

\item Its Gram matrix is
\[
Q:=MM^{\mathsf T}
  =A^{3}+2A^{2}+2I,
\]
where \(A\) is the adjacency matrix of \(\Gfifteen\).

\item The adjacency spectrum is
\[
\spec(A)
=
\left\{
4^{(1)},\,
2^{(5)},\,
(-1)^{(4)},\,
(-2)^{(5)}
\right\},
\]
and therefore
\[
\spec(Q)
=
\left\{
98^{(1)},\,
18^{(5)},\,
3^{(4)},\,
2^{(5)}
\right\}.
\]

\item The signed lift
\[
\Gthirty\longrightarrow\Gfifteen
\]
contains nontrivial holonomy information not determined by \(Q\).

\item Declared projections in the higher finite construction can
return before the corresponding lifted or complete registered state
returns.

\item Registered history denotes retained finite information that
distinguishes complete states already identified by a declared
projection.

\item A declared real-valued register
\(h\in\mathbb R^n\) has centered relational magnitude
\[
R_n(h)
=
\left\|
h-\overline h\,\one_n
\right\|_2.
\]
This scalar is a many-to-one compression, not a complete state
coordinate.

\item The native registered carrier has \(60\) directed-edge
contextual states.  Every directed quotient edge has exactly two
registered continuations, one of designation \(b\) and one of
designation \(ab\), and every live oriented two-edge path has a unique
registered completion.

\item The accepted registered language contains \(120\) oriented
words.  Its exact completion profile is
\[
8\longrightarrow2\longrightarrow1\longrightarrow1\longrightarrow1.
\]

\item The five-address registration partitions the fifteen-state
sector register into a five-dimensional fibre-constant sector and a
ten-dimensional fibre-internal sector.

\end{enumerate}

These structures remain distinct.  In particular,
\[
\text{native selection}
\neq
\text{quadratic aggregation}
\neq
\text{covering holonomy}
\neq
\text{registered history}
\neq
\text{centered scalar magnitude}.
\]

\subsection{New contribution}

The upgraded theorem package has four parts.

First, for a live registered prefix \(p\), let
\[
\mathcal C(p)
\]
denote the lawful completion family.  Prefix extension is monotone:
\[
p\preceq p'
\quad\Longrightarrow\quad
\mathcal C(p')\subseteq\mathcal C(p).
\]
The certified native loading ladder is
\[
8\longrightarrow2\longrightarrow1.
\]
We therefore define native relational loading as accumulated
admissible constraint and registered-future saturation by
\[
|\mathcal C(p)|=1.
\]

Second, future uniqueness does not erase complete context.  The
two-edge saturated stage contains \(120\) distinct contexts with
\(120\) distinct registered completions, and the terminal completion
register retains two native holonomy designations over every closing
chamber.

Third, let \(\Delta\) denote the context-activated native
\(b/ab\) designation response on the fifteen-state chamber register.
Let \(H\) denote the exact ten-dimensional five-address-hidden
projector.  Then
\[
\Delta^{\mathsf T}\Delta
=
\Delta\Delta^{\mathsf T}
=
\frac34H.
\]
The range and domain of \(\Delta\) lie in the hidden sector, while
the five-address readout satisfies
\[
\rho\Delta=0.
\]

Fourth, the Thalean quadratic provides a separate
contrast-concentrating operation.  For
\[
\widehat Q:=\frac1{98}Q,
\]
the common mode is fixed and the complete centered subspace contracts
by the sharp operator factor \(9/49\).  Consequently,
\[
\widehat Q^{\,k}x
\longrightarrow
\Pzero x.
\]
Every finite iterate remains invertible.

The native loading response and the centered Thalean quadratic lie in
the same exact three-dimensional transport-response algebra, but the
two operators are not proportional and do not perform the same
operation.

\subsection{Native loading is not quadratic contrast concentration}

The distinction can be tested directly using the Paper 14 contrast
observable.

For the canonical uniform completion-occupancy register, native
loading gives
\[
\zeta:
\frac{31}{55}
\longrightarrow
\frac{41}{65}
\longrightarrow
\frac{7}{10}.
\]
The common component remains fixed while contrast strictly increases
on every one of the \(180\) strict loading extensions.

Thus
\[
\text{native relational loading}
\neq
\text{\(\widehat Q\)-type contrast concentration}.
\]

Native loading sharpens the selected future.

The normalized quadratic suppresses centered contrast relative to
common mode.

The theorem keeps these operations separate.

\subsection{Theorem boundary}

The following implications are not asserted:
\[
Q
=
\text{native temporal evolution},
\]
\[
\widehat Q
=
\text{native loading},
\]
\[
\widehat Q
=
\text{physical averaging law},
\]
or
\[
\text{relational saturation}
=
\text{physical gravitational collapse}.
\]

The exact theorem is finite and internal.

The native carrier, the retained-sector response, and the quadratic
contrast law are joined mathematically without identifying any of
them with a physical process.

\subsection{Three epistemic layers}

The paper maintains three layers.

\begin{description}

\item[Established finite mathematics.]
The graph, quotient hierarchy, incidence matrix, quadratic identity,
spectra, native carrier language, completion-family loading order,
context activation, projection-relative saturation, retained hidden
sector, covering receipts, recurrence distinctions, registered
history, Thalean time, and centered magnitude.

\item[Internal conditional mechanics.]
Clean release, overcompression, any proposed composition of native
selection with contrast concentration into a larger cycle, and any
stronger global inter-face transport construction.

\item[Physical bridge.]
Any proposed relation to physical information density, gravity,
black holes, horizons, singularities, the CMB, or Planck-scale
support.

\end{description}

No result is promoted from one layer to the next without an explicit
crosswalk.
